\selectlanguage{ngerman}%
\frenchspacing%
\begin{ledgroupsized}[r]{120mm}%
\footnotesize%
\pstart%
\noindent\textbf{Überlieferung:}
\pend%
\end{ledgroupsized}
\begin{ledgroupsized}[r]{114mm}%
\footnotesize%
\pstart%
\parindent -6mm%
\makebox[6mm][l]{\textit{L}}%
Notiz:
LH~XXXV~11,\;14 Bl.\;32.
Ein Zettel (etwa 9\:x\:5~cm);
unterer Rand beschnitten.
Eine Seite auf Bl.\;32\:r\textsuperscript{o}; Bl.\;32\:v\textsuperscript{o} ist unbeschrieben.
\pend%
\end{ledgroupsized}%
\vspace{5mm}%
\begin{ledgroup}%
\footnotesize%
\pstart%
\noindent%
\textbf{Datierungsgründe:}
Die undatierte Notiz N.\;\ref{58947} ist auf einem Träger überliefert, der weder Spüren eines Wasserzeichens aufweist noch sonstige datierungsrelevante Anhaltspunkte bietet.
Für eine grobe chronologische Einordnung erweisen sich folgende inhaltliche Beobachtungen als ausschlaggebend:
\pend%
\pstart%
(1)~%
Die Notiz N.\;\ref{58947} beginnt mit der Feststellung, die Kinematik des schiefen sowie des mehrfachen Stoßes lasse sich durch das Modell eines \glqq Schiffes im Schiffe\grqq\ adäquat beschreiben.
(Hierunter ist eine Erweiterung der "Schiffsmethode" gemeint, d.\,h. jener Methode der Zerlegung zusammengesetzter Bewegungen, die in Leibnizens stoßtheoretischen Untersuchungen seit spätestens Mitte 1677 Anwendung fand.)
Bei den genannten Stoßarten falle dieses Modell %
zudem mit der in N.\;\ref{58947} nicht weiter erläuterten \textit{methodus quantitatis potentiae respectivae per diversas compositiones conservandae} zusammen (\refpassage{LH_35_11_14_032_methodus-1}{LH_35_11_14_032_methodus-2}).
Mit dem exzentrischen Stoß verhalte es sich aber \textendash\ wie der Schlussatz in N.\;\ref{58947} einräumt \textendash\ anders:
Das Modell eines \glqq Schiffes im Schiffe\grqq\ lasse sich hierauf nicht anwenden, obwohl die \glqq Natur des exzentrischen Stoßes\grqq\ durch Leibnizens eigene Methode (\textit{mea methodus}) doch erfolgreich ermittelt worden sei. %
Ist bei dieser Anspielung nun wieder die \glqq Methode der bei verschiedenen Zusammensetzungen zu bewahrenden relativen Potenz\grqq\ gemeint, so ist hieraus zu schließen, dass diese Methode aus Leibnizens Sicht nicht ohne Weiteres mit dem \glqq Modell eines Schiffes im Schiffe\grqq\ zusammenfällt, da ihr offenbar ein breiterer Anwendungsbereich zukommt als letzterem, dem sich zumindest der exzentrische Stoß entzieht. 
\pend%
\pstart%
(2)~%
Die \glqq Methode der bei verschiedenen Zusammensetzungen zu bewahrenden relativen Potenz\grqq\ wird in N.\;\ref{58947} ganz unvermittelt genannt.
Daher ist wahrscheinlich, dass die Notiz später entstand als die Texte, die diese Methode im Zusammenhang mit dem exzentrischen Stoß erläutern bzw. ausführen:
Das sind u.\;a.
die vermutlich im Herbst\;(?) 1690 verfasste Aufzeichnung N.\;\ref{58949}, in der die besagte Methode konzise dargelegt wird (\refpassage{LH_35_11_14_034_034rv_drstllng-1}{LH_35_11_14_034_034rv_drstllng-2} im vorliegenden Band),
sowie
die eigenhändig auf den 31.\;Oktober\,/\,10.\;November 1690 datierte \textit{Scheda secunda de impulsu eccentrico} N.\;\ref{58939}, in der die Methode ausführlich angewendet wird (\refpassage{LH_35_11_14_003-004_004v_methodus-1}{LH_35_11_14_003-004_004v_methodus-2} im vorliegenden Band).
Da Leibniz im Schlusssatz von N.\;\ref{58947} die erfolgreiche Anwendung eben dieser Methode auf den exzentrischen Stoß erwähnt, ist anzunehmen, dass die Notiz %
nach dem Entwurf N.\;\ref{58939} und wohl auch nach der Aufzeichnung N.\;\ref{58949} verfasst wurde.
\pend%
\pstart%
(3)~%
Es gilt ferner zu bemerken, dass die \glqq Methode der bei verschiedenen Zusammensetzungen zu bewahrenden relativen Potenz\grqq\ auch in der auf den Zeitraum 1690\textendash1694 datierbaren und in diesem Band edierten Aufzeichnung N.\;\ref{60298} besprochen wird, dort im Hinblick auf den Zusammenstoß dreier Körper und ohne Berücksichtigung des exzentrischen Stoßes. %
Auch in N.\;\ref{60298} wird die besagte Methode ausdrücklich mit dem Modell eines \glqq Schiffes im Schiffe\grqq\ verglichen, welchem ebenfalls mit Skepsis begegnet wird. %
Somit lässt sich auch in der Aufzeichnung N.\;\ref{60298} eine Diskrepanz zwischen der allgemeinen Methode und dem fiktionalen Modell feststellen.
Die Notiz N.\;\ref{58947} zeigt indessen, dass die Zweifel am erweiterten Schiffsmodell im Fall des multiplen (und somit schiefen) Stoßes überwunden wurden, nicht aber im Fall des exzentrischen Stoßes.
Dies dürfte belegen, dass N.\;\ref{58947} nach N.\;\ref{60298} entstand.
\pend%

\pstart%
(4)~%
In Untersuchungen über den exzentrischen Stoß wird die einfache \glqq Schiffsmethode\grqq\ mehrfach in Anspruch genommen, ohne grundsätzliche Zweifel an deren Angemessenheit zu erheben.
Dies ist etwa bei folgenden, sämtlich im vorliegenden Band edierten Texten der Fall:
der Aufzeichnung \textit{De impulsu eccentrico} N.\;\ref{60300}, \refpassage{LH_37_05_044_044r_navis-1}{LH_37_05_044_044r_navis-2} (April--Oktorber 1690);
der \textit{Scheda de impulsu eccentrico} N.\;\ref{58940}, \refpassage{LH_35_11_14_005-006_005v_navis-1}{LH_35_11_14_005-006_005v_navis-2} (31.\;Oktober\,/\,10.\;November 1690);
den zusammenhängenden Texten \textit{Casus eccentrici concursus simplicissimus} N.\;\ref{60106}, Randbemerkungen zu \refnote{FnA:35_10_13_007r_navis} u. \refnote{FnA:35_10_13_007v_navis} und N.\;\ref{60107}, \refpassage{LH_35_10_13_008_008rv_navisficta-1}{LH_35_10_13_008_008rv_navisficta-2} (September 1691).
Hieraus würde sich -- unter der Annahme, dass die Zweifel am Modell des "Schiffes im Schiffe" eine weitere, spätere Entwicklung darstellen -- \textit{September 1691} als frühester Entstehungszeitpunkt der Notiz N.\;\ref{58947} ergeben.
\pend%
\pstart%
(5)~%
Die Notiz N.\;\ref{58947} könnte allerdings durch ihre Zweifel am Modell \glqq des Schiffes im Schiffe\grqq\ auch mit dem im vorliegenden Band edierten Entwurf N.\;\ref{41780} verbunden sein, in welchem ähnliche Zweifel an derselben Methode erhoben werden.
Der auf die Zeitspanne von Mitte September 1690 bis 1696 datierbare Entwurf N.\;\ref{41780} zeigt somit, dass die Annahme einer späteren Entwicklung, aus der die Zweifel an der Erweiterung der Methode sich ergeben hätten, fraglich ist und dass die Notiz N.\;\ref{58947} %
bereits im Laufe des Jahres 1690 abgefasst worden sein könnte.
Eine Entstehung vor den Texten N.\;\ref{58939} und N.\;\ref{58949} sowie N.\;\ref{60298} ist jedoch -- wie oben festgehalten wurde -- unwahrscheinlich.
Hieraus ergibt sich als mutmaßlicher frühester Entstehungszeitpunkt der Notiz N.\;\ref{58947} der \textit{10.\;November 1690}.
Eine spätere Datierung -- ab September 1691 -- ist aber nicht unwahrscheinlich.
\pend%
\pstart%
(6)~%
Ein eindeutiger Terminus ante quem der Datierung lässt sich für die Notiz N.\;\ref{58947} zum Redaktionsschluss nicht angeben.
\pend%
\end{ledgroup}%
\selectlanguage{latin}%
\frenchspacing%
\vspace{8mm}%
\normalsize%
\pstart%
\noindent%
\edtext{%
\lbrack%
\href{\Scan{LH_35_11_14_032r}}%
{32\,r\textsuperscript{o}}%
\rbrack\ %
Casus%
\indexS{casus solvendus}
concursus obliqui%
\indexS{concursus obliquus}
et compositi%
\indexS{concursus compositus}%
}{%
\lemma{%
\href{\Scan{LH_35_11_14_032r}}{\lbrack32\:r\textsuperscript{o}\rbrack}%
}\Bfootnote{%
\textit{(1)}~Casus unius corporis in plura incurrentis variique alii ob%
\indexS{concursus plurimus}\indexS{casus varius}
\textit{(2)}~Casus concursus obliqui et compositi%
~\textit{L}}}
solvi possunt
supposita navi in%
\indexS{navis in navi}
\edtext{%
navi.
Quo idem consequemur,
quod mea
\edlabel{LH_35_11_14_032_methodus-1}%
methodo%
\indexS{methodus mea}%
\indexS{methodus potentiae respectivae conservandae}%
\indexS{quantitas potentiae respectivae}%
\indexS{potentia respectiva conservanda}
quantitatis potentiae respectivae
per diversas compositiones%
\indexS{compositio motuum}
conservandae.%
\edlabel{LH_35_11_14_032_methodus-2} %
\edlabel{LH_35_11_14_032_natura-0}%
Sed explicari non poterit
(\protect\vphantom)%
quantum hactenus judicare possum%
\protect\vphantom()
\edlabel{LH_35_11_14_032_natura-1}%
natura concursus%
\indexS{natura concursus eccentrici}%
\indexS{concursus eccentricus}%
}{%
\lemma{navi.}\Bfootnote{%
\textit{(1)}~Sed ad concursus
\textit{(2)}~Quo idem %
\lbrack...\rbrack\ conservandae. Sed
\textit{(a)}~solvi
\textit{(b)}~explicari non poterit
\textit{(aa)}~casus
\textit{(bb)}~(\protect\vphantom)quantum hactenus %
\lbrack...\rbrack\ natura concursus%
~\textit{L}}}
eccentrici
\edtext{%
quem mea methodus%
\indexS{methodus mea}
\edtext{etiam}{\lemma{etiam}\Bfootnote{\hspace{0.5mm}%
\textit{erg.\;L}}}
absolvit.%
\edlabel{LH_35_11_14_032_natura-2}%
}{%
\lemma{quem \lbrack...\rbrack\ absolvit}\Cfootnote{%
Siehe hierüber die ausführlichen Datierungsgründe.
}}
\pend%
\selectlanguage{latin}%
\frenchspacing%
\normalsize%
\newpage%
